\chapter*{Introduction générale}
\addcontentsline{toc}{chapter}{Introduction générale}
\markboth{Introduction générale}{Introduction générale}
\label{chap:introduction}
%\minitoc

Le secteur événementiel connaît une évolution importante, portée par la digitalisation des services liés à la découverte, à la réservation et à la gestion des événements. Les participants recherchent aujourd’hui des solutions simples leur permettant de consulter les événements disponibles, d’obtenir des informations détaillées, de réserver leurs billets, d’effectuer leurs paiements en ligne et d’accéder rapidement à leurs tickets. De leur côté, les organisateurs ont besoin d’outils efficaces pour créer leurs événements, gérer leur billetterie, suivre leur activité et organiser le contrôle d’accès.

Cependant, les solutions existantes restent souvent fragmentées. Certaines plateformes se limitent à la promotion des événements, tandis que d’autres se concentrent sur la billetterie, le paiement ou le contrôle d’accès. Cette dispersion complique le parcours des utilisateurs et rend plus difficile le travail des organisateurs, notamment pour le suivi des ventes, la revente des tickets, le scan à l’entrée et la supervision des opérations sensibles.

Dans ce contexte, la mise en place d’une plateforme numérique centralisée, intitulée FESTY, représente une réponse pertinente aux besoins actuels du secteur événementiel. La plateforme vise à regrouper, dans un même environnement, les principales fonctionnalités nécessaires aux visiteurs, aux utilisateurs, aux partenaires organisateurs, aux agents de scan et aux administrateurs. Elle comprend une application mobile, un espace partenaire et un back-office d’administration.

Dans le cadre de notre projet de fin d’études, nous avons effectué un stage au sein de TEKSIGHT, une entreprise spécialisée dans les services informatiques et les solutions numériques. Notre mission consiste à concevoir et développer la plateforme FESTY afin de faciliter la découverte des événements, la personnalisation de l’expérience utilisateur, la gestion des partenaires, la billetterie, le paiement, la revente des tickets, le contrôle d’accès et la supervision administrative.

Ce rapport présente les différentes étapes de réalisation du projet, depuis l’étude du besoin jusqu’à la conception et la réalisation des fonctionnalités principales. Il met également en évidence les choix méthodologiques, fonctionnels et techniques adoptés, ainsi que l’organisation du projet selon une démarche Scrum découpée en quatre sprints.

Le premier chapitre présente le contexte général du projet, l’organisme d’accueil TEKSIGHT, la problématique, l’étude de l’existant, la solution proposée, la méthodologie de travail adoptée et la planification prévisionnelle du projet. Le deuxième chapitre est consacré à l’analyse et à la préparation du projet. Il présente les acteurs, les besoins fonctionnels et non fonctionnels, le backlog produit, la planification des sprints, les diagrammes globaux, l’environnement de travail, les technologies utilisées et l’architecture générale de la solution.

Les chapitres suivants présentent l’étude et la réalisation des quatre sprints du projet. Le troisième chapitre traite le Sprint 1, consacré à l’authentification, à l’inscription, à la vérification du numéro de téléphone, à la gestion du profil utilisateur et à l’accès sécurisé au back-office administrateur. Le quatrième chapitre aborde le Sprint 2, dédié à la gestion des partenaires et des événements, depuis la demande partenaire jusqu’au contrôle administratif de l’offre événementielle.

Le cinquième chapitre présente le Sprint 3, centré sur l’exploration des événements et des stars, la personnalisation, les recommandations, les favoris, le chat événementiel, la gestion du line-up artistique et le catalogue artistique. Enfin, le sixième chapitre est consacré au Sprint 4, qui couvre la billetterie, le paiement, la revente, la consultation des tickets, le contrôle d’accès par scan, le suivi partenaire et la supervision des opérations financières sensibles.
